\section{Introduction}

Cardinal characteristics of the continuum are cardinals between $\aleph_1$,
the first uncountable cardinal, and $\mathfrak c$, the cardinality of the
continuum. They measure combinatorial properties of the real numbers; see
\cite{blass2009combinatorial} for background.

The symbol $\omega$ denotes the set of natural numbers, which includes $0$ and is identified with the first infinite ordinal.

The cardinal characteristics $\mathfrak{ap}$ and $\mathfrak{dp}$, both related
to the pseudointersection number $\mathfrak p$, were studied by Brendle in
\cite{brendle1999dow}. We recall the relevant definitions.

\begin{itemize}
    \item For a set $S$, $[S]^{<\omega}$ denotes the collection of all finite subsets of $S$, and $[S]^\omega$ denotes the collection of all countably infinite subsets of $S$.
    \item Sets $a,b\in[\omega]^{\aleph_0}$ are orthogonal, written $a\perp b$, if $a\cap b$ is finite.
    \item Collections $\mathcal A,\mathcal B\subseteq[\omega]^{\aleph_0}$ are orthogonal, written $\mathcal A\perp\mathcal B$, if $a\perp b$ for every $a\in\mathcal A$ and $b\in\mathcal B$.
    \item If $\mathcal A,\mathcal B\subseteq[\omega]^{\aleph_0}$, then $\mathcal A$ is weakly separated from $\mathcal B$ if there exists $X\subseteq\omega$ such that $a\not\perp X$ for every $a\in\mathcal A$, and $b\perp X$ for every $b\in\mathcal B$. In this case, we say that $X$ weakly separates $(\mathcal A,\mathcal B)$.
    \item $\mathcal A\subseteq[\omega]^{\aleph_0}$ is an almost disjoint family if $a\perp b$ for every two distinct $a,b\in\mathcal A$.
\end{itemize}

Our convention regarding the order of weak separation follows Banakh and Bazylevych \cite{banakh2023q}.  It is the reverse of Brendle's convention for ordered pairs \cite{brendle1999dow}.

Since weak separation is not symmetric, the order of the two families matters in
the following definitions.

\begin{itemize}
    \item $\mathfrak{dp}$ is the least cardinal $\kappa$ for which there exist $\mathcal A,\mathcal B\subseteq[\omega]^{\aleph_0}$ such that $\mathcal A\perp\mathcal B$, $|\mathcal A\cup\mathcal B|=\kappa$, and $\mathcal A$ is not weakly separated from $\mathcal B$.
    \item $\mathfrak{ap}$ is the least cardinal $\kappa$ for which there exist $\mathcal A,\mathcal B\subseteq[\omega]^{\aleph_0}$ such that $\mathcal A\perp\mathcal B$, $|\mathcal A\cup\mathcal B|=\kappa$, $\mathcal A\cup\mathcal B$ is an almost disjoint family, and $\mathcal A$ is not weakly separated from $\mathcal B$.
\end{itemize}

It is immediate that $\mathfrak{dp}\leq\mathfrak{ap}$.
This inequality can be strict: Brendle showed by forcing that
$\mathfrak{dp}<\mathfrak{ap}$ is consistent \cite{brendle1999dow}.
Both cardinals are greater than or equal to $\mathfrak p$, and consistently
strictly greater than $\mathfrak p$ \cite{brendle1999dow}. Here $\mathfrak p$ is the least
cardinality of a centered family of infinite subsets of $\omega$ with no
pseudointersection; that is, every nonempty finite subfamily has infinite
intersection, but there is no single infinite set $P$ such that
$P\setminus A$ is finite for every member $A$ of the family.

If $X$ is a topological space and $Y\subseteq X$, we say that $Y$ is a $G_\delta$ subset of $X$ if there are open subsets $(U_n)_{n \in \omega}$ of $X$ such that $Y=\bigcap_{n<\omega} U_n$.
Dually, $Y$ is an $F_\sigma$ subset of $X$ if it is a countable union of closed subsets of $X$.
We say that $X$ has a $G_\delta$-diagonal if the set $\Delta_X=\{(x,x):x\in X\}$ is a $G_\delta$ as a subset of $X^2$.

A $Q$-set is classically defined as an uncountable set of reals all of whose subsets are
$G_\delta$ sets in the subspace topology (equivalently, all of its subsets are
$F_\sigma$ sets in the subspace topology) \cite{miller1984special,fleissner1980qsets}. The existence of $Q$-sets is independent of the Zermelo--Fraenkel set theory with the axiom of choice (ZFC) \cite{fleissner1980qsets,miller1984special}. Indeed, if
there is a $Q$-set, then it has at least $2^{\aleph_1}$ subsets, all realized as
traces of $G_\delta$ subsets of the reals. Since there are only $\mathfrak c$
many $G_\delta$ subsets of the reals, no $Q$-set can exist under the Continuum
Hypothesis (CH). On the other hand, every set of reals of cardinality less than
$\mathfrak{ap}$ has all of its subsets $G_\delta$ in the subspace topology
\cite{brendle1999dow}; hence, if
$\aleph_1<\mathfrak{ap}=\mathfrak c$, as happens for example under Martin's
Axiom together with the negation of CH \cite{blass2009combinatorial}, then $Q$-sets
exist.

$Q$-sets are also related to the separable case of the Normal Moore Space
Conjecture: a non-metrizable separable normal Moore space exists if and only if
there is a $Q$-set \cite{heath1964screenability}.

The smallest cardinality of a non-$Q$-set is denoted by $\mathfrak q$
\cite{miller1984special,brendle1999dow}.
More generally, a $Q$-space is a topological space in which every subset is a
$G_\delta$ set. In this terminology, $\mathfrak q$ is the least cardinality of
a metrizable separable non-$Q$-space, equivalently, of a regular second-countable
non-$Q$-space.

These notions have been studied in more general settings. There are two natural
directions to pursue: weaken regularity to weaker separation axioms, or weaken
second countability. In \cite{banakh2023q}, Banakh and Bazylevych followed the
first direction and defined the cardinals $\mathfrak q_1$, $\mathfrak q_2$ and
$\mathfrak q_{2\frac{1}{2}}$ by considering second-countable non-$Q$-spaces
which are, respectively, $T_1$, Hausdorff and Urysohn. Recall that a topological
space is $T_1$ if every singleton is closed; Hausdorff, or $T_2$, if every two
distinct points have disjoint neighborhoods; and Urysohn, or
$T_{2\frac{1}{2}}$, if every two distinct points have disjoint closed
neighborhoods.

Banakh and Bazylevych also introduced a cardinal $\mathfrak{adp}$, which lies
between $\mathfrak{dp}$ and $\mathfrak{ap}$ and is useful in the study of
$Q$-spaces.

\begin{definition}
    $\mathfrak{adp}$ is the least cardinal $\kappa$ for which there exist $\mathcal A,\mathcal B\subseteq[\omega]^{\aleph_0}$ such that $\mathcal A\perp\mathcal B$, $|\mathcal A\cup\mathcal B|=\kappa$, $\mathcal A$ is an almost disjoint family, and $\mathcal A$ is not weakly separated from $\mathcal B$.
\end{definition}

They asked the following.

\begin{prob}[{\cite[Problem 12]{banakh2023q}}] Is $\mathfrak{adp}=\mathfrak{dp}$? Is $\mathfrak{adp}=\mathfrak{ap}$?
\end{prob}

In Section~\ref{sec:adp}, we answer the first part by proving
$\mathfrak{adp}=\mathfrak{dp}$ in ZFC. Together with the consistency of
$\mathfrak{dp}<\mathfrak{ap}$, this shows that
$\mathfrak{adp}=\mathfrak{ap}$ is not provable in ZFC.

The asymmetry of weak separation also suggests asking what happens if the
almost disjointness requirement in the definition of $\mathfrak{adp}$ is imposed
on $\mathcal B$ instead of on $\mathcal A$.

\begin{definition}
    $\mathfrak{adp}_2$ is the least cardinal $\kappa$ for which there exist $\mathcal A,\mathcal B\subseteq[\omega]^{\aleph_0}$ such that $\mathcal A\perp\mathcal B$, $|\mathcal A\cup\mathcal B|=\kappa$, $\mathcal B$ is an almost disjoint family, and $\mathcal A$ is not weakly separated from $\mathcal B$.
\end{definition}

Using the same technique, we show that $\mathfrak{adp}_2=\mathfrak{ap}$.
Thus, in every model in which $\mathfrak{dp}<\mathfrak{ap}$, this modified
cardinal is distinct from $\mathfrak{dp}$.

Banakh and Bazylevych also asked:

\begin{prob}[{\cite[Problem 11(1)]{banakh2023q}}] Is $\mathfrak{ap}\leq\mathfrak q_2$?
\end{prob}

In Section~\ref{sec:hausdorff-q-spaces}, we provide a positive answer.
We introduce a tree analogue of
$\mathfrak{ap}$, denoted by $\mathfrak{at}$. This cardinal is obtained by
replacing almost disjoint families of subsets of $\omega$ with almost
disjoint families of infinite subtrees of an $\omega$-tree. It satisfies
$\mathfrak{ap}\leq\mathfrak{at}$, and we show that $\mathfrak{at}=\mathfrak q_2$ (Theorem~\ref{thm:at-equals-q2}), yielding $\mathfrak{ap}\leq\mathfrak q_2$.

In Section~\ref{sec:small-diagonals}, we define the variants
$\mathfrak q_{i\Delta}$ obtained by additionally requiring a
$G_\delta$-diagonal. We prove
$\mathfrak q_{1\Delta}=\mathfrak{at}=\mathfrak q_2$ and
$\mathfrak q_{2\frac{1}{2}\Delta}=\mathfrak q_{2\frac{1}{2}}$.

Finally, in Section~\ref{sec:ap-below-at-model}, we show that the tree cardinal
$\mathfrak{at}$ is not simply another presentation of $\mathfrak{ap}$.
Assuming the Generalized Continuum Hypothesis (GCH), for every regular cardinal $\lambda>\omega_1$ we construct a ccc
forcing extension in which
\[
    \mathfrak{ap}=\omega_1<
    \mathfrak{at}=\mathfrak q_{2}=\mathfrak c=\lambda.
\]
In particular, $\mathfrak{ap}<\mathfrak q_2$ is consistent.
The construction is based on Brendle's model for $\mathfrak{ap}<\mathfrak q$.

\subsection*{Summary of results}

We end the introduction with a summary of the main results of this paper.

\begin{itemize}
    \item We prove $\mathfrak{adp}=\mathfrak{dp}$ in ZFC (Theorem~\ref{thm:adp-equals-dp}), answering the first question in \cite[Problem~12]{banakh2023q}. We define a dual variant, $\mathfrak{adp}_2$, and we prove $\mathfrak{adp}_2=\mathfrak{ap}$ (Theorem~\ref{thm:adp2}).

    \item We introduce the tree cardinal $\mathfrak{at}$ and prove $\mathfrak{ap}\leq\mathfrak{at}=\mathfrak q_2$ (Theorem~\ref{thm:at-equals-q2}). This gives a combinatorial characterization of $\mathfrak q_2$ and answers \cite[Problem~11(1)]{banakh2023q} positively: $\mathfrak{ap}\leq\mathfrak q_2$ holds in ZFC.

    \item We determine the effect of requiring a $G_\delta$-diagonal in the definition of $Q$-space cardinals. We prove $\mathfrak q_{1\Delta}=\mathfrak{at}=\mathfrak q_2$ (Theorem~\ref{thm:q1delta-equals-at}). We also show that every second-countable Urysohn space has a $G_\delta$-diagonal, yielding $\mathfrak q_{2\frac{1}{2}\Delta}=\mathfrak q_{2\frac{1}{2}}$ (Lemma~\ref{lem:second-countable-urysohn-gdelta-diagonal}).

    \item We show that the inequality $\mathfrak{ap}\leq\mathfrak q_2$ is consistently strict. Assuming GCH, for every regular $\lambda>\omega_1$ we construct a ccc forcing extension with $\mathfrak{ap}=\omega_1<\mathfrak{at}=\mathfrak q_2=\mathfrak q_{2\frac{1}{2}}=\mathfrak c=\lambda$ (Theorem~\ref{thm:consistent-ap-below-at}). Thus the inequality answering \cite[Problem~11(1)]{banakh2023q} cannot be strengthened to an equality in ZFC.
\end{itemize}
